% 本文件由 LaTeX 排版 Agent 根据有效 translation.json 自动生成。
% author=Augustine of Hippo; work_id=1311; title=论守寡的益处

\chapter{论守寡的益处}

% structure: p0001 source=h1 target=omitted reason=page-title-already-rendered-by-parent

\OriginalTerm{论守寡的益处}{De bono viduitatis}\par

本作品在\WorkTitle{订正录}中未见提及，很可能因为它是一封信函，故\Person{波西迪乌斯}将其列入第七章。它也被标记为\Person{圣奥古斯丁}之作，因其引用其其他作品，如\WorkTitle{论婚姻的益处}等第十五章，以及\WorkTitle{致普罗巴书}第二十三章。日期以\Person{德米特里亚斯}最近被祝圣为标记，即413年。\Person{犹利亚纳}在\WorkTitle{书信188}中感谢他给予的劝诫，可能是指反对伯拉纠主义的劝诫。\par

有人以其与398年\WorkTitle{第四次迦太基大公会议}第104条不符为由提出反对意见，该会议将祝圣后再婚的寡妇处以绝罚，并宣告她们犯奸淫；而在第十章和第十一章中，反对此类婚姻无效且应回归守贞的观点被驳斥。然而，两者并非完全不可调和，因为可能犯有类似奸淫的罪过，并以绝罚形式施加谴责，同时婚姻仍然有效。\WorkTitle{迦克墩大公会议}第16条对此规定了刑罚，主教有权减轻。——\Emph{节选自本笃会版}\par

\Person{奥古斯丁}主教，基督的仆役，也是基督仆役们的仆役，致天主虔敬的婢女\Person{犹利亚纳}：在主中之主内获享健康。\par

1. 为不再欠你因基督而发出的请求与爱所许下的承诺，我趁着机会，在我其他极为紧迫的事务中，写信给你，略论神圣守寡的益处；因为当我与你同在时，你以恳求加于我，而我未能拒绝你之后，你常以书信要求我履行承诺。在这部作品中，当你阅读时若发现有些事完全不涉及你本人，或与你同在基督内生活的人无关，且对你的生活并非严格必需，你切不可因此认为它们是多余的。因为此信虽写给你，却不可只为你一人而写；但我们确实不可忽视，它当借你之手有益于他人。因此，你在此发现的一切，无论是否曾为或今为你所需，且你觉察到对他人有益者，当乐于拥有并借阅，使你的爱德也对他人生益。\par

2. 因此，在一切关于生活与行为的问題上，不仅需要教导，也需要劝勉；以便我们通过教导知道该做什么，通过劝勉被激励，不认为去做已知当做之事是烦难的。我能教导你什么，超过我们读自宗徒的？因为圣经为我们的教导立下标准，使我们不敢\BibleQuote{聪明超过所当聪明的}；而是如他自己所说的，要\BibleQuote{按照天主分给各人信德的大小，清醒地}聪明。因此，我不可教导你任何其他事，只当向你解释那位导师的话，并照主所赐予我的来论述它们。\par

3. 因此，外邦人的导师、蒙选之器宗徒说：\BibleQuote{但我对未婚者和寡妇说：他们若常像我一样，为他们是有益的。}这些话当这样理解，即我们不可认为寡妇不应被称为未婚者，因为她们似乎已体验过婚姻；因为他用\OriginalTerm{未婚女子}{unmarried woman}一词指现在不被婚姻约束的妇女，无论她们曾否结婚。他在另一处讲明了这一点，说：\BibleQuote{未婚女子和童贞女是分开的。}显然，当他加上\Quote{童贞女}时，他所谓的\OriginalTerm{未婚女子}{unmarried woman}除了寡妇还能是什么？因此，在接下来的经文中，他用\OriginalTerm{未婚者}{unmarried}这同一术语涵盖两种身份，说：\BibleQuote{未婚的女子挂虑主的事，怎样使主喜悦；但已婚的女子挂虑世上的事，怎样使丈夫喜悦。}无疑，他理解的\OriginalTerm{未婚}{unmarried}不仅指从未结婚者，也指因守寡而从婚姻束缚中释放、不再结婚者；因为他也因此不称\Quote{已婚}者，除非她有丈夫，而不包括曾有丈夫而今无者。因此，每个寡妇都是未婚者；但因并非每个未婚者都是寡妇（也有童贞女），所以他在此同时提到了二者，说：\BibleQuote{但我对未婚者和寡妇说：}就好像他说：我对未婚者所说的，不仅对童贞女说，也对寡妇说：\BibleQuote{他们若常像我一样，为他们是有益的。}\par

4. 请看，你的善与宗徒称为他自己的那善相比，倘若信德存在；是的，更确切地说，因为信德存在。这教导简短，但不可因此被轻视，因为它简短；反而更易于牢记，也更宝贵，因为在简短中它并不廉价。因为宗徒在此所要陈明的并非任何种类的善，他明确地将其置于已婚妇女的信德之上。然而，已婚妇女的信德——即那些在婚姻中结合的基督徒和虔诚妇女的信德——有多么伟大，可以从这一点理解：当他为躲避淫乱而给予训诫时（他无疑也在对已婚者说话），他说：\BibleQuote{你们不知道你们的身体是基督的肢体吗？} 因此，信实婚姻的善如此之大，以至于连肢体也是（肢体）基督的。但是，既然守寡的节制之善比这善更好，那么这一奉献的目的并不是要使一位天主教寡妇比基督的肢体更多，而是要使她在基督的肢体中比已婚妇女有更好的位置。因为同一宗徒说：\BibleQuote{就如我们在一个身体上有许多肢体，但肢体并非都有同一作用；同样，我们众人在基督内也是一个身体，彼此都是肢体，按着赐给我们的恩宠，各有不同的恩赐。}\par

5. 为此，当他也劝告已婚者不要彼此剥夺肉体的应得之份时，免得因此，其中一方（婚姻之债被拒绝后）因自己的不节制受撒殚诱惑而堕入淫乱，他说：\BibleQuote{这是我说，出于宽免，并非出于命令；但我愿意众人如同我自己一样；但每人都有自己从天主得来的恩赐；有人这样，有人那样。} 你看，婚配的贞洁和基督徒床笫的婚姻信德，也是一种恩赐，且出自天主；因此，当肉欲略微超出感官交合所需的尺度，超出生育子女所必需时，这恶并非来自婚姻，而是因婚姻之善而可获宽恕的小罪。因为关于婚姻——为生育子女而缔结，以及婚配贞洁的信德，和婚姻的圣事（只要双方都活着就不可解除）——这些都是善；但关于那种在两性之事上不节制的用法，这在已婚者的软弱中被认可，并因婚姻之善的介入而得宽恕，宗徒说：\BibleQuote{我说这话是出于宽免，并非出于命令。} 同时，当他说：\BibleQuote{女人在丈夫活着的时候是被束缚的；但如果她的丈夫死了，她就自由了：她可以随意嫁人，只要在主内；但如果她照我的劝告继续守寡，她就更有福了。} 他充分表明，一个忠信的女子在主内是有福的，即使她在丈夫死后第二次结婚；但寡妇在同一位主内更有福。也就是说，不仅用圣经的话，也用圣经的例子来说，\Person{卢德}是有福的，但\Person{亚纳}更有福。\par

6. 因此，首先你应当知道，你所选择的善并不谴责第二次婚姻，而是将其置于较低的地位。因为，正如你的女儿所选择的至圣童贞之善并不谴责你的第一次婚姻；同样，你的守寡也不谴责任何人的第二次婚姻。因为，正是由此，\OriginalTerm{卡特弗里吉亚派}{Cataphryges}和\OriginalTerm{诺瓦替安派}{Novatians}的异端膨胀，\OriginalTerm{戴尔都良}{Tertullian}也鼓着充满声音而非智慧的双颊，用辱骂的牙齿攻击第二次婚姻，仿佛是非法的；而宗徒以清醒的头脑允许完全合法。从这健全的教义，不要让任何人的推理——无论他无知还是有学识——动摇你；也不应如此赞扬你自己的善，以致指责他人的善为恶（其实并非恶）；但你应因自己的善更加喜乐，因为你越看到，借此不仅避开了恶，而且超越了一些善。因为奸淫和淫乱是恶。但远离这些非法之事的人，她借着誓愿的自由约束了自己，并且不是出于法律的命令，而是出于爱德的劝告，使得即使是合法之事对她也不再合法。婚姻的贞洁是善，但守寡的节制是更好的善。因此，这更好的善因那另一善的服从而得尊荣，并非那另一善因对这更好之善的赞美而被定罪。\par

7. 但是，既然宗徒在称赞独身男女的果实，因为他们思念主的事，怎样悦乐天主时，接着又说：\BibleQuote{我说这话是为你们的益处，并不是要给你们设下圈套，}也就是说，并不是要强迫你们；\BibleQuote{而是为那体面的事；}我们就不应因为他称独身者的善为体面，就以为婚姻的束缚是卑鄙的；否则我们也要谴责初婚了，连\OriginalTerm{弗吕家派}{Cataphryges}、\OriginalTerm{诺瓦替安派}{Novatians}，以及他们最有学问的拥护者\OriginalTerm{戴尔都良}{Tertullian}都不敢说这是卑鄙的。但是，当他这样说：\BibleQuote{我对未婚者和寡妇说：如果他们能像我一样继续守独身，对他们是有益的；}他确实用\Quote{好}来代替\Quote{更好}，因为每一件与一善比较而被称作更好的事，无疑也是善；因为除了更为善之外，它还能被称作更好的什么呢？然而我们并不因此就推断他认为他们若结婚就是恶，因为他曾说：\BibleQuote{如果他们能继续这样，对他们是有益的；}同样，当他说\BibleQuote{而是为那体面的事}时，他并没有表明婚姻是卑鄙的，而是将那比（另一件）体面的事更体面的事，用体面的一般名称来称赞。因为什么是更体面？不就是更体面吗？但更体面的一定是体面的。事实上他清楚地表明这比那件善事更好，当他说：\BibleQuote{谁嫁女，做得好；但不嫁女，做得更好。}并且这比那有福的更有福，当他说：\BibleQuote{但如果她继续这样，将更有福。}因此，正如有比善更好的，比有福更有福的，同样也有比体面更体面的，他选择称之为体面。因为远非如此，那被\Person{伯多禄}宗徒所说的话是卑鄙的：\BibleQuote{丈夫们，要以敬重待妻子，如同对待较软弱的器皿，与她们一同承受恩宠的生命，}又向妻子们讲话，以\Person{撒辣}为榜样劝她们服从自己的丈夫：\BibleQuote{因为从前那些仰望天主的圣善妇人，也是这样装饰自己，顺从自己的丈夫；就如\Person{撒辣}服从了\Person{亚巴郎}，称他为主；你们若是行善，不畏惧任何惊扰，就成为了她的女儿。}\par

8. 因此，宗徒\Person{保禄}论及未婚女子时所说的话：\BibleQuote{使她身心成圣，}我们不可理解为，一位已婚而贞洁、并按照圣经顺从丈夫的信女，在身体上不成圣，只在精神上成圣。因为当灵魂被圣化时，身体也不可能不成圣，圣化的灵魂使用身体；但为了避免让人觉得我们只是在争论而没有用神圣的言语证明这一点；既然\Person{伯多禄}宗徒提到\Person{撒辣}时，只说\BibleQuote{圣善妇人}，而没有说\BibleQuote{及在身体上}；让我们思想同一\Person{保禄}的另一句话，在禁止奸淫时他说：\BibleQuote{你们不知道你们的身体是基督的肢体吗？那么，我可以将基督的肢体作为娼妓的肢体吗？决不可！}所以，有谁敢说基督的肢体是不圣洁的呢？或谁敢将有信仰已婚者的身体从基督的肢体中分离出去呢？因此，稍后他又说：\BibleQuote{你们的身体是圣神的宫殿，这圣神是你们从天主领受的，你们不是属于自己的人，因为你们是重价买来的。}他说信友的身体既是基督的肢体，又是圣神的宫殿，这其中无疑包括男女两性的信友。因此那里有已婚女子，也有未婚女子；但她们的功绩不同，如同肢体相对于肢体更优越，然而没有一个肢体从身体分离。因此，当他论及未婚女子说\BibleQuote{使她身心成圣}时，他要人理解为在身心两方面更圆满的圣化，并没有剥夺已婚女子身体的所有圣化。\par

9. 因此，要明白你的善，不，更要记住你已经学到的，你的善更受赞扬，是因为有另一件善，比这一件更好，而不是因为若没有那一件，这一件就不能成其为善，除非那一件是恶，或根本不存在。眼睛在身体中大有尊荣，但如果只有眼睛，没有其他较不尊荣的肢体，眼睛的尊荣就会减少。在天上，太阳以其光辉超越月亮，并不是责备月亮；此星与彼星的荣光不同，并不是因骄傲而相争。因此，\BibleQuote{天主造了一切，看，都很好；}不仅是\Quote{好}，而且是\Quote{甚好}；没有别的原因，只是因为\Quote{一切}。因为对每一件个别工程也说过：\BibleQuote{天主看了认为是好的。}但是，当\Quote{一切}被提及时，就加上了\Quote{甚好}；并且说：\BibleQuote{天主看了他所造的一切，看，都很好。}因为某些个别事物比另外一些个别事物更好，但全部一起比任何个别更好。因此，愿基督纯正的教义藉着祂的恩宠使你在祂的身体内得以健全，使你那统治身体的灵魂，对于你在身心方面比别人所拥有的更好，既不傲慢地自夸，也不愚昧地分别。\par

10. 然而，我称\Person{卢德}为有福，\Person{亚纳}更为有福，乃因前者两次结婚，后者则在丧夫守寡后长久生活，莫非你就立刻以为自己比\Person{卢德}更好？诚然，在先知时代，圣善妇人的处置是不同的：是服从而非情欲迫使她们结婚，为的是天主子民的繁衍，以使基督的先知在他们中被预先派遣；而子民本身，借着在他们中以预像发生的事——无论是那些知道的人，还是那些不知道的人——也无非是基督的先知，基督的肉身也应由他而生。因此，为了那个子民的繁衍，凡在以色列不留下后裔的，便按法律的规定被视为可咒骂的。圣善的妇女也被激动，不是出于肉体交合的情欲，而是出于生育的虔诚；所以我们极有理由相信，倘若家庭能以其他方式延续，她们决不会寻求肉体交合。并且，丈夫被允许同时拥有多位妻子；而其原因并非肉欲，而是对生育的远见，这由以下事实表明：正如圣善的男人可以有多位妻子在世，圣善的女人却不得与多位在世丈夫交合；因为她们越是追求那不能增加其多产的事，就越显卑下。因此，圣\Person{卢德}因没有当时以色列所必需的后嗣，在丈夫死后便寻求另一人以便得子。因此，相较这位两次结婚者，一次结婚即守寡的\Person{亚纳}更为有福，因为她还得以成为基督的女先知；关于她，我们应当相信，尽管她未有儿子（圣经虽缄默不语，对此未作确证），但她藉那圣神预见到基督即将由童贞女诞生，并得以认出他还是孩童；因此，即使无子（即假定她无子嗣），她也有充分理由拒绝第二次婚姻：因她知道，现在是更好地事奉基督的时候，不是借着生育的义务，而是借着节制的热忱；不是借着已婚子宫的多产，而是借着寡妇生活的贞洁。但若\Person{卢德}也知道，通过她的肉身繁衍了后裔，基督日后将由这后裔取得肉身，并且她以结婚来表明对此认知的服务，我就不敢再断言\Person{亚纳}的寡居比她多产更为有福。\par

11. 但是你，既有儿子，又活在这世界的末期，如今不是抛掷石头的时节，而是聚敛的时节；不是拥抱的时节，而是禁戒拥抱的时节；当时宗徒大声疾呼说：\BibleQuote{弟兄们，我对你们说：时日不多了；从此以后，有妻子的要像没有一样；}的确，若你寻求第二次婚姻，那将不再是先知或法律的服从，甚至也不是对子嗣的肉欲，而只是不节制的记号。因为你所做的，正是宗徒在说过\BibleQuote{为他们更好，如果他们能像我一样继续守贞}之后所说的；他立刻接着又说：\BibleQuote{但若他们不能节制，就让他们结婚；因为我宁愿他们结婚，胜过被欲火焚烧。}他说这话，为的是使不受约束的欲望之恶不至被婚姻的尊贵状态所接纳而一头栽入罪恶的卑劣中。但感谢主，因你生下了你所不愿成为的，而你孩子的童贞弥补了你童贞的失落。因为基督信仰的教义，经仔细询问后答称：如今这时节，就连第一次婚姻也应被轻视，除非不节制阻挡在前。因为那说\BibleQuote{若他们不能节制，就让他们结婚}的人，本可以说\Quote{若他们没有儿子，就让他们结婚}，如果现在，在基督复活与宣讲之后，有如此众多且丰盛的属灵儿子可生于万民之中，那么按肉身生儿子仍像初时一样是一种义务的话。然而，他在另一处说：\BibleQuote{但我愿意年轻的寡妇结婚，生育儿女，管理家务}，他以宗徒的节制与权威称赞了婚姻之善，但并没有将生育的义务，如同遵守法律一样，强加给那些\Quote{接受}节制之善的人。最后，他为何说这话，在他接着又说的话中显明了：\BibleQuote{免得给敌人毁谤的机会；因为有些人已经转去随从了撒殚}；通过这些话语我们可理解，那些他愿意他们结婚的人，本可以节制而不结婚，但结婚比随从撒殚更好，即比回头向后看那已过去且消逝的事物，从而从童贞或寡居的贞洁崇高目的中堕落更好。因此，凡不能节制的，让他们在宣示节制之前、在向天主许愿之前结婚；若他们不还愿，就会被公正定罪。诚然，在另一处他论及这样的人说：\BibleQuote{因为他们既在基督内过着安逸的生活，就想结婚；他们被定罪，因为废弃了起初的信德；}即他们将自己的意志从节制的目的转向了婚姻。诚然，他们废弃了那信德，藉此他们先前曾许下不愿恒心持守的愿。因此，婚姻之善确实常是善：但在天主子民中，它曾是服从法律的行为；如今它是软弱的补救，但在某些人身上是人性的慰藉。诚然，致力于得子——不是像狗一样滥交雌性，而是借着婚姻的正当秩序——并不是我们可以责备之人的情感；然而，这情感本身，思念天上之事的基督徒心智，以更可赞的方式超越和克服了它。\par

12. 但是，既然正如主所说：\Quote{不是所有人都能接受这话；}因此，能接受的，就接受吧；不能节制的，就结婚吧；尚未开始的，要三思；已经承担了的，要坚持；不可给仇敌以机会，不可使祭献从基督那里被夺走。诚然，在婚姻的纽带中，若保存了贞洁，就不必害怕定罪；但在寡居和童贞的节制中，所追求的是一种更大恩赐的卓越；而且，当这恩赐已被追求、选择，并以誓愿的债务奉献之后，从那时起，不但进入婚姻，甚至虽然没有结婚而只是愿意结婚，就成了被定罪的事。因为，为了表明这一点，宗徒并没有说：\Quote{当她们在享乐中生活时，在基督内}结婚；而是说：\Quote{她们愿意结婚；因此，}他说，\Quote{她们有定罪，因为她们废弃了起初的信德。}虽然不是因结婚而废弃，乃是因愿意而废弃；并非连这样的婚姻也被判定为定罪之事；但受定罪的，是善意所受的损害，是誓愿之信德的破坏，受定罪的，不是因较低之善而得的缓解，而是从更高之善的堕落：最后，这样的人被定罪，并非因她们后来进入了婚姻的信德，而是因她们废弃了起初节制的信德。并且，为了用简短的话提示这一点，宗徒并没有说那些在立定更大圣善的目的之后结婚的人有定罪（不是因为他们不被定罪，而是免得在他们身上婚姻本身被认为该定罪）；而是在说了\Quote{她们愿意结婚}之后，立即加上\Quote{她们有定罪}。他又说明了理由：\Quote{因为她们废弃了从前的信德，}为要表明，是那从目的堕落的意愿被定罪，无论结婚与否。\par

13. 因此，那些说这种婚姻不是婚姻，而是通奸的人，依我看，并没有足够敏锐和细心地考虑他们所说的话；诚然，他们被一种真理的表象所误导。因为，那些因基督徒的圣洁而不结婚的人，据说是选择了与基督的婚姻，由此某些人争辩说：如果她在丈夫活着的时候嫁给了别人，她就是通奸者，正如主自己在福音中所规定的；因此，在基督活着的时候（死亡不能再管辖祂），如果那选择了祂婚姻的女子嫁给了一个人，她就是通奸者。说这话的人，的确敏锐，但却没有注意到这推论事实上带来了多么大的荒谬。因为，即使在丈夫活着的时候，经他同意，一个女子向基督誓愿守贞，本是值得称赞的；但按照这些人的推理，没有人应该这样做，免得她使基督自己（这是不敬的想象）成为一个通奸者，因为她在丈夫活着的时候与祂结婚。其次，尽管初婚比再婚有更好的功德，但圣洁的寡妇们切不可有这样的想法，即基督对她们来说像是第二个丈夫。因为她们从前也拥有祂自己（当她们顺服并忠信地服事自己的丈夫时），不是按照肉体，而是按照圣神作丈夫；祂是教会（她们是其肢体）的丈夫；教会因信德、望德、爱德的正直，不仅在童贞女中，而且在寡妇和忠信的有夫之妇中，全然是一个贞女。诚然，宗徒对普世教会（她们都是其肢体）说：\Quote{我把你们许配给一个丈夫，如同一个贞洁的童女，献给基督。}但是，祂知道如何使一个妻子成为贞女，在不损贞洁的情况下结果实，祂自己的母亲甚至在肉体上也能不损贞洁地怀上祂。但由于这种欠考虑的观念（他们因此认为那些从这圣洁目的堕落的女子若结婚，就不是婚姻），带来了不小的恶，即妻子们与她们的丈夫分离，仿佛她们是通奸者而不是妻子；而那些希望使这样分离的女子恢复节制的想法，却使她们的丈夫成为真正的通奸者，因为他们在妻子活着的时候娶了别人。\par

14. 因此，对于从更好的目的堕落的妇女，如果她们结了婚，我实在不能说她们是通奸而不是婚姻；但我明确毫不迟疑地说，从更神圣的贞洁（这是向主誓愿的）中的偏离和堕落，比通奸更糟糕。因为，如果（这是绝不可怀疑的）基督的一个肢体对她的丈夫不忠，这属于对基督的冒犯；那么，当对祂自己（在祂所要求献上的事上）不忠时，这冒犯又该是多么严重呢？而祂本来并没有要求这事应该被献上。因为，当一个人未能归还他所誓愿的（不是出于命令的强制，而是出于劝告的建议）时，他违反誓愿的不义，因他当初誓愿的必要性越小，就增加得越多。我讨论这些事，是为了使你不会认为再婚是犯罪的，或者认为任何婚姻（就其是婚姻而言）是一种恶。因此，你应当有这样的心思：不是要谴责她们，而是要轻视她们。因此，寡居贞洁之善以一种更光辉的方式是适宜的，因为为了发愿和宣认它，妇女可以轻视那既愉快又合法的事。但在发愿宣认之后，她们必须继续驾驭和克服那愉快的事，因为它不再合法了。\par

15. 人们往往提出关于第三次或第四次婚姻，甚至比这更多次婚姻的问题。对此要严格作答，我不敢谴责任何婚姻，也不敢夺去这些多次婚姻的羞耻。但是，唯恐我这简短的回答偶然使任何人不满，我准备听我的驳斥者更详尽地论述。因为他或许举出某种理由，说明为什么第二次婚姻不被谴责，而第三次却被谴责。因为我，如同在这篇讲论的开始我所警告的，不敢比应当有的智慧更有智慧。因为我是谁，竟认为必须定义我认为宗徒未曾定义的事呢？因为他说：\BibleQuote{女人活着的时候，是受束缚的，只要她的丈夫活着。}他没有说：第一个，或第二个，或第三个，或第四个；而是说：\BibleQuote{女人}，他说，\BibleQuote{活着的时候是受束缚的，只要她的丈夫活着；但如果她的丈夫死了，她就自由了；她可以随意嫁人，只要在主内；但如果她能这样坚持下去，她就更有福了。}我不知道关于此事，在这句话上还能添加或删除什么。其次，我听见宗徒们和我们自己的师尊与主亲自回答撒杜塞人，当时他们向他提出一个女人，不是结过一次婚或两次婚，而是，如果可以这样说，结过七次婚，在复活时她是谁的妻子？他斥责他们说：\BibleQuote{你们错了，既不了解圣经，也不明了天主的能力。因为在复活的时候，他们不娶也不嫁，因为他们不再会死，而要像天主的使者一样。}因此他提到了那些复活进入生命的人，而不是那些复活受罚的人。因此他本可以说：你们错了，不了解圣经，也不明了天主的能力，因为在那复活中，不可能有曾是多人之妻的女子；然后又补充说，那里也没有人结婚。但是，正如我们所看到的，他在这句话中也没有显示任何谴责那位有如此多丈夫的妻子的迹象。因此，我也不敢违背天然的羞愧感，说：当她的丈夫们死后，一个女人可以随意嫁多少次；我也不敢出于自己的内心，在圣经的权威之外，谴责任何数量的婚姻。但是，我对一个只有一个丈夫的寡妇所说的话，我对每一个寡妇都说：如果你能这样坚持下去，你就更有福了。\par

16. 因为那个通常被提出的问题也不是愚蠢的：谁能说哪一位寡妇在功绩上更应被拣选？是那位只有一个丈夫，在与其丈夫共同生活了相当长的时间后，成了寡妇，有儿子出生并活着，且宣示了守寡的节操；还是那位作为年轻女子，在两年内失去了两个丈夫，没有活着的孩子可以安慰她，却向天主誓许了守寡，并在其中以最持久的圣德变老？在此，让他们讨论，如果可能，向我们展示一些证据，我们来衡量寡妇的功绩，不是根据丈夫的数目，而是根据节操本身的力量。因为，如果他们说过，那位只有一个丈夫的应被优先于那位有两个丈夫的，除非他们举出某种特殊的理由或权威，他们必定会被发现是看重灵魂的优点，但不是更伟大的灵魂优点，而是肉体的好运。因为与丈夫长期生活并怀孕生子，都属于肉体的好运。但是，如果他们并非因为这个原因（她有儿子）而偏爱她，那么她和丈夫长期生活的事实本身，除了是肉体的好运之外，还能是什么呢？再者，\Person{亚纳}自己的功绩在此主要受到称赞，因为她如此早地埋葬了丈夫之后，通过她漫长的生命，长久地与肉体斗争并得胜了。因为经上这样记载：\BibleQuote{有一位女先知\Person{亚纳}，是\Person{法努尔}的女儿，属\Place{阿协尔}支派。她年纪已经老迈，从作童女出嫁后，同丈夫住了七年，就寡居了，直到八十四岁，总不离开\Place{圣殿}，禁食祈祷，昼夜事奉天主。}你看，这位圣洁的寡妇不仅在这一点上受称赞：她只有一个丈夫，而且她从童女时代与丈夫只生活了短短几年，并以如此虔诚的服务，将寡妇的贞洁职守坚持到如此高寿。\par

17. 因此，让我们把三位寡妇放在眼前，每一位拥有那全部东西中的一项：我们假设一位只有一个丈夫，但她缺少如此长的寡居时间（因为她与丈夫生活了很长时间），也缺少如此大的虔诚热忱（因为她没有那样以禁食和祈祷事奉）；第二位，在前夫非常短的生命之后，很快又失去了第二个丈夫，并且现在长期寡居，但她自己也没有那样致力于最虔诚的禁食祈祷事奉；第三位，她不仅有两个丈夫，而且还与每个丈夫（或其中一位）长期生活，在生命后期成了寡妇，诚然，如果她愿意结婚，她也可以怀孕生子，但她接受了寡居的节操；但她更专注于天主，更小心地常常做祂喜欢的事，日夜如同\Person{亚纳}一样，以祈祷和禁食事奉。如果问题被提出，其中哪一位在功绩上更应被拣选，谁能不看到在这场竞赛中，棕榈枝必须给予那更伟大、更炽热的虔诚呢？同样，如果另外三位被设定，每一位拥有那三项中的两项，但各缺其中一项，谁能怀疑他们会是更好的，他们将在他们的两样好处中拥有更卓越的虔诚谦卑，以便有崇高的虔诚呢？\par

的确，这六位寡妇中没有一位能达到你的标准。因为如果你将这一誓愿持守到老年，你就可能拥有\Person{安娜}所擅长的一切三样美善。因为你曾有过一位丈夫，且他与你血肉同居的时间不长；这样，如果你顺从宗徒的话，即\BibleQuote{但真正寡居且孤獨的婦人，已盼望於主，日夜堅持祈禱}，并以清醒的警惕避开随后所说的\BibleQuote{但那貪圖享樂的寡婦，雖活猶死}，那么\Person{安娜}所拥有的一切三样美善，也将属于你。但你也生了儿子，而她可能没有。然而，你并非因此而被称赞，不是因为你有了他们，而是因为你热切地以虔诚养育和教育他们。因为他们为你所生，是出于多产；他们活着，是出于好运；他们被如此抚养，是出于你的意愿和安排。在前者，让别人恭喜你；在这件事上，让别人效法你。\Person{安娜}通过先知性的认识，认出了基督与祂的童贞母亲；而福音的恩宠却使你成为基督的一位童贞女的母亲。因此，那位圣童贞女，她自己愿意并寻求，你就把她献给了基督，这为她的祖母和母亲的寡居功绩增添了童贞的功绩。因为拥有她的你，不会缺少从她而来的某种东西；在她身上，你成为了你自身所不是的。因为圣童贞女在你结婚时被夺去，正是为此事发生，为要使她由你而生。\par

因此，这些关于已婚妇女不同功绩和不同寡妇不同功绩的讨论，我本不会在这部著作中开始，如果我写给你的是仅仅写给你一个人的话。但是，由于这类论述中有某些非常困难的问题，我希望说一些比与你直接相关更深入的话，因为有一些人，他们除非试图不只是通过判断来讨论，而是通过撕碎来切烂别人的辛勤劳动，否则就不觉得自己有学问；其次，也为了你本人不仅可以持守你所誓愿的，并在这善事上进步，而且更仔细、更确定地知道，你的这一善并非与婚姻的恶区分开来，而是在婚姻的善之前。因为不要让那些谴责寡妇再婚的人——尽管他们通过戒除许多你所使用的事物而操练节欲——因此迷惑你，使你认同他们的想法，尽管你无法做他们所做的事。因为没有人会因此成为疯子，尽管他看到疯子的力气比神志健全的人更大。因此，首要的是，让纯正的道理既装饰又守护志善。事实上，正是因为这个原因，天主教妇女，即使她们曾结婚不止一次，在公义的判断中，不仅优于仅有一位丈夫的寡妇，也优于异端者的童贞女。的确，关于婚姻、寡居和童贞这三件事，有许多迂回曲折的问题和许多困惑；要通过讨论深入探究并解决这些，需要更仔细的关注和更充分的论述，使我们在所有这一切中或是持有正确的心意，或是在任何事上持有别样的心意，愿天主也将这启示给我们。然而，宗徒随后也说了：\BibleQuote{我們已經到達何處，就在何處行走。} 但是，就我们所谈论的这件事而言，我们已经达到了把节欲置于婚姻之前，而把圣童贞甚至置于寡居的节欲之前；并且不因我们自己或朋友任何志趣的赞美而谴责任何婚姻，只要这些婚姻不是奸淫而是婚姻。关于这些事，我们在一部名为\WorkTitle{论婚姻的好处}的书中，在另一部名为\WorkTitle{论圣童贞}的书中，以及在一部我们竭尽全力撰写的\WorkTitle{驳摩尼教徒福斯特}的书中，已经说了许多其他事；因为他在著作中以最尖刻的责难攻击圣祖和先知的贞洁婚姻，使某些无知者的心偏离了信德的纯正。\par

因此，既然在这篇小作品的开头，我提出了两件必要的事，并承担了逐一论述：一件与教义有关，另一件与劝勉有关；在前一部分，我已尽己所能，按所承担的任务没有懈怠；现在我们进入劝勉部分，好使那被明智认识的美善能被热切追求。在这件事上，我首先劝告你：无论你感到自己心中对虔诚守寡的爱有多大，你都要将其归功于天主的恩宠，并向祂谢恩——祂藉其圣神白白赐予你如此之多，以致当祂的爱倾注在你心中时，对更美好之事的爱将合法之事的许可从你那里夺去。因为这是祂的恩赐：当结婚合法时，你就不愿结婚；而如今，即使你愿意，它也已不合法；这样，不愿意行此事的心愿便更为坚定，以免现在不合法的事被行出，而它即使合法时也未曾行出；并且，作为基督的寡妇，你竟能达到如此地步，得以看见你的女儿也成为基督的贞女；因为当你像亚纳那样祈祷时，她已成为玛利亚所是的样子。你越知道这些是天主的恩赐，你就越因这些恩赐而蒙福；事实上，你之所以如此，无非是因为你知道从谁那里获得了你所拥有的。因为请听宗徒对此事所说的话：\BibleQuote{但我们所领受的，不是这世界的精神，而是由天主而来的圣神，为使我们知道天主所赐予我们的事。}诚然，许多人拥有天主的许多恩赐，却因不知道从谁那里获得，而以不敬的虚妄自夸。但没有人因天主的恩赐而蒙福，却对施与者忘恩负义。此外，在神圣礼仪中，我们被召叫\Quote{举心向上}，我们之所以能够如此，是藉着祂的帮助——我们因祂的诫命而受劝勉；因此，对于这举心向上的巨大美善，我们不归光荣于自己，好似出自自身的力量，而是向我们的主天主谢恩。因为我们立刻被提醒：\Quote{这是合宜的}，\Quote{这是正义的}。你记得这些话出自何处，你认识它们被何种权威、以何等大的圣德在内里被赞美。因此，持守并保有你已领受的，并向施与者谢恩。因为，虽然领受并拥有是你的本分，但你所有的，是你所领受的；因为对于那些骄傲自大、不敬地夸耀自己所有、仿佛出于自身的人，主通过宗徒说：\BibleQuote{你有什么不是领受的呢？既然是领受的，为什么你夸耀，好像不是领受的呢？}\par

21. 我不得不因某些人的小言论而加以劝诫，这些言论应当躲避和远离，它们已开始潜入许多人的耳中，进入他们的思想，这些言论（说来令人流泪）敌视 \OriginalTerm{基督的恩宠}{grace of Christ}，力图劝服我们，认为我们不必向主祈求，使我们不陷于诱惑。因为他们如此试图维护人的自由意志，以为单凭自由意志，甚至无需天主恩宠的帮助，我们就能履行天主所命令的事。由此得出，主曾徒然说：\BibleQuote{你们应醒寤祈祷，免得陷于诱惑；}而我们在主祷文中每日也徒然念：\BibleQuote{不要让我们陷于诱惑。}因为若不凭我们自己力量就不能胜过诱惑，为何我们祈求不进入，也不被引入诱惑呢？我们倒不如行我们所自由意愿、且最绝对之能事；并嘲笑宗徒所说：\BibleQuote{天主是忠信的，他决不许你们受那超过你们能力的试探；}且反对他，说：为何我向主寻求那已置于我自己能力之内的事？但心思正直的人绝不如此思想。因此，让我们寻求祂赐予祂命令我们拥有之物。为此祂命令我们拥有这尚未拥有之物，为劝诫我们该寻求什么；并且当我们获得履行祂所命令的能力时，我们也就能明白，这也是从谁那里领受的；免得我们被这世界的精神所膨胀、高举，而不知道天主赐予我们什么。为此，我们绝不破坏\OriginalTerm{人的自由意志}{free choice of the human will}，当我们不因忘恩的骄傲而否认，反而因感恩的虔诚而彰显那扶助这自由意志本身的天主恩宠时。因为意愿属于我们：但意愿本身被劝醒而兴起，被治愈而得能力，被扩展而接受，被充满而拥有。如果我们不愿意，当然既不会领受所赐之物，也不会拥有。因为谁会有\OriginalTerm{节德}{continence}（在论述天主其他恩赐时，尤其我要对你说的这个），谁，我说，会有节德，除非自愿？因为也没有人会自愿接受。但若你问，这是谁的恩赐，使它能被我们的意愿接受并拥有？请听圣经；是的，更好说，因为你认识，回忆你所读过的：他说：\BibleQuote{我知道，除非天主赐与，谁也不能节德，而这也是智慧，以知恩赐来自何处。}这两个恩赐伟大：\OriginalTerm{智慧}{wisdom}与节德；智慧，诚然，使我们按天主的认识被塑造；节德，则使我们不与此世同化。但天主命令我们既要有智慧又要有节德，没有这些美善，我们不能成为正义和成全的。然而，让我们祈求祂赐予祂所命令的，借着帮助和感动，祂已借着命令和召唤劝诫我们该意愿什么。凡祂已赐予的这些，让我们祈求祂保全；但尚未赐予的，让我们祈求祂补充；同时，让我们为我们所领受的感恩并祈求；至于尚未领受的，正因我们对自己所领受的不忘恩，我们当信我们将获得。因为祂已赐予已婚的信徒能力，使他们在奸淫和私通上节制，祂自己也将能力赐予神圣的童贞者和寡妇，使她们在一切性交上节制；对于这一德行，如今恰当地使用\OriginalTerm{无玷的贞洁}{inviolate chastity}或节德一词。或者，莫非我们确实从祂领受了节德，而智慧却来自我们自己？那么，雅各伯宗徒所说的是什么：\BibleQuote{你们中谁若缺乏智慧，就该向那慷慨赐与而不责斥的天主祈求，必会赐给他。}但关于这个问题，在我们其他的小作品中，在主帮助我们的范围内，我们已经说了许多；而在其他时候，借着祂所能赐予的，当有机会时，我们将继续谈论。\par

22. 如今，我因我们某些极其友善亲爱的弟兄而有意对此主题说些话；他们虽无故意之罪，却确实陷于这错误之中，但仍陷于其中；这些人认为，当他们劝勉他人归向正义与虔敬时，其劝勉若要有力，就必须将他们意欲在人身上促成之事（即人应行之事）完全置于人的能力之内，而非仰赖天主的恩宠相助，仅凭自由意志独自发出；仿佛人能凭自由意志行善，除非藉天主的恩赐获得释放！他们没有留意到，他们自己之所以拥有如此劝勉的能力——能激动人迟钝的意志去开始善生、点燃冷漠、纠正迷失、挽回偏离、平息对抗——这本身也是天主的恩赐。因为这样他们就能成功说服他们所要说服的事；否则，如果他们不在人的意志中成就这些事，那么他们的工作又是什么？他们为何还要说话？不如任人随其自择。但若他们确实在人的意志中成就了这些事，那么请问：人凭言语在人的意志中成就如此大事，难道天主在那里毫无帮助吗？不，无论人的言词多么有力，能藉讨论的技巧和言语的甘美，在人的意志中植入真理、滋养爱德、藉教导消除错误、藉劝勉消除懈怠，但\Quote{栽种的不算什么，浇灌的也不算什么，唯有那使之生长的天主，才算得什么}。因为工人若不在外面用尽一切方法，而造物主不在里面隐秘地工作，一切都是枉然。因此，我盼望这封信藉您贤德的善举能尽早送到这些人手中；为此，我以为应当就这主题说些话。其次，为使您本人，以及任何阅读或听人诵读此信的寡妇，知道你们在爱慕和实践守贞之善上，藉自己的祈祷比藉我们的劝勉进步更大；因为若我们的言论对你们也有助益，这一切必须归功于祂的恩宠，正如经上所载：\Quote{祂的手中有我们和我们的言论。}\par

23. 因此，如果你们尚未向天主誓许寡妇的守贞，我们定会劝你们誓许；但既然你们已经誓许，我们就劝你们恒心持守。然而，我看也必须这样说话，好引导那些至今仍考虑结婚的人也爱慕并把握这善。因此，让我们听从宗徒的话：\Quote{未出嫁的妇女，挂虑主的事，一心使身体和灵魂圣洁；但已出嫁的妇女，挂虑世俗的事，想怎样悦乐丈夫。}他并非说，挂虑世俗的事，以至于不能圣洁；而是肯定说，婚姻的圣洁较小，因为有一部分挂虑是关乎取悦世俗的快乐。因此，凡一个热切心志的女子会用在取悦丈夫之事上的精力，未婚的基督徒妇女理应某种程度地收集起来，用于那取悦主的努力。请想一想，她取悦的是谁——取悦主的人，她越取悦祂，就越有福；但她越多思虑世俗的事，就越少取悦祂。所以，你们要以全部的热忱取悦祂，祂是\Quote{超越人子之美}的那一位。因为你们能取悦祂，是由于祂的恩宠，即\Quote{倾注在祂唇上的}。请你们也把那原会为取悦丈夫而占据心神的思虑部分，用来取悦祂。请取悦祂——祂曾不取悦世界，好使那些取悦祂的人从世界获得释放。因为这位超越人子之美的，人们在十字架的苦难中看见祂，\Quote{祂没有俊美，也没有华彩，祂的面容憔悴，形貌丑陋。}然而，从你们救赎主的这丑陋中，流出了你们美丽的代价——但这是内心的美丽，因为\Quote{君王女儿的美丽全在内里}。你们要凭这美丽取悦祂，要以勤勉和焦虑的思念来装饰这美丽。祂不爱虚假的染料；真理喜爱真实的事，而祂——若你们认出所读的——被称为真理。祂说：\Quote{我是道路、真理、生命。}你们要藉着祂奔向祂，凭祂取悦祂；与祂同在、在祂内、因祂而生活。要以真实的情感和最圣洁的贞洁，爱被这样一位新郎所爱。\par

24. 愿那童贞女，即你圣洁的孩子，也用内心的耳朵听到这些。我将看她在那位君王的国度中在你面前能走多远：那是另一个问题。然而，你，母亲和女儿，已找到了那一位，你们当凭贞洁之美一同取悦祂；她轻视了婚姻，而你则轻视了第二次婚姻。诚然，如果有需要你们取悦的丈夫，或许此时你们会羞于和女儿一同打扮自己；如今，不要让羞愧阻止你们共同致力于那能使你俩一同增光的事；因为你们一同被那一位所爱，不是指责，而是荣耀。但那些用脂粉伪造的红白之色，即使有丈夫，你们也不会使用；因为你们不认为他们是值得欺骗的人，也不认为自己是应当欺骗的人。如今，那位君王，祂渴慕祂唯一净配（你们是她的肢体）的美貌，你们当以完全的真诚一同取悦祂，一同依附于祂：她以童贞的贞洁，你以寡妇的节制，两人都以属灵的美貌。在这美貌中，她的祖母，你的婆婆，此时想来已经年迈，也与你们一同美丽。的确，当爱德将这美貌的活力带向未来的事物时，岁月的流逝不会在它上面留下皱纹。你们有一位圣洁的老年妇人，同在你们的家中，也在基督内，你们可向她请教关于恒心的事：你们该如何与这种或那种诱惑作战，该做什么才能更容易地克服它；该采取什么防护措施，使诱惑不再轻易设伏；以及诸如此类的事。她教导你们，她因时间而稳固，因爱德而心怀善意，因亲情而充满忧虑，因年龄而无忧无虑。你们要特别向她请教，她已体验过你们所体验的。因为你的孩子所唱的那首歌，在 \BibleBook{若望}{默示录} 中只有童贞女才能唱。但她为你们二人祈祷，比为自己更谨慎；她更为她的孙女忧虑，因为孙女有更长的岁月需要克服诱惑；而你，她看到更接近自己的年龄，且是这样一个年龄的女儿的母亲，若你看到她结婚（如今这是不合法的，但愿她远离此事），我想你会羞于与她一同生育子女。那么，如今危险年龄还剩下多少呢？你之所以不被称为祖母，正是为了与女儿一同在圣善思想和行为的后裔上结果实。因此，祖母更为她忧虑，并非没有理由，因为她的誓愿更大，而她刚刚开始的一切仍在她前面。愿主俯听她的祈祷，使你圣洁地追随她的功德：她在青年时生下了你丈夫的肉身，在老年时又为你的女儿心灵受产痛。所以，你们众人，同心合意，以品行取悦，以祈祷催促那一位独一净配的唯一丈夫，在祂的肢体内你们由同一圣神而生活。\par

25. 已逝之日不再返回，今日紧随昨日，明日又紧随今日；看，所有的时间和世间事物都在逝去，为叫那应许的永恒得以来临；\BibleQuote{谁若坚持到底，这人才可得救。} 如果世界如今正在腐朽，那结婚的妇人，她为谁生育？或是心中将要生育，而肉体不再生育，她为何结婚？但如果世界仍将存续，为何不更爱那创造世界的主？若今生的诱惑已经消逝，基督徒灵魂便没有值得渴望追求的事物；但若它们仍存留，便有值得以圣德轻视之物。因为这两者之一，没有情欲希望，另一则有更大的爱德光荣。甚至那些肉性年龄似乎开花结果的岁月，又有多少或多长呢？有些女子怀着结婚的念头，并热切渴望，但因被轻视或拖延，突然老去，以致如今她们更以结婚为羞，而非渴望。但许多结婚的女子，丈夫婚后不久便远行，她们在等待中老去，好似很快成为寡妇，有时甚至未能至少在老年时迎接丈夫归来。因此，若已订婚的新郎被轻视或延误，或丈夫在外时，肉欲能克制而不犯奸淫或通奸，为何不能克制而不犯亵圣之罪？若它在被延迟时炽热而能被抑制，为何在断绝后冷却反而不能克服？因为那些不绝望于同一种快感的人，更强烈地忍受着欲望的炽热。但那些向天主誓守贞洁的未婚者，恰恰抽走了那希望，即爱的燃料。因此，更容易约束那因无期待而点燃的欲望；然而，若不为此祈祷以克服它，它作为非法之物反而更被热烈渴望。\par

26. 因此，在圣洁的守寡生活中，应以属灵的喜乐取代肉身的喜乐：阅读、祈祷、圣咏、善念、多行善工、对来世的希望，以及向上的心；并为这一切感谢光明的圣父，毫无疑问，一切美好的恩赐和完美的恩赐，正如圣经所证明的，都从他那里降下。\newline{}ewline{}因为当已婚妇女在她们丈夫肉身中的喜乐被其他肉身的喜乐所取代，仿佛是为了安慰她们时，我又何必谈论随之而来的邪恶呢？因为宗徒已简洁地说过：\Quote{那生活在享乐中的寡妇，虽生犹死}。\newline{}ewline{}但愿你绝不至于被财富的贪欲所占据，以取代婚姻的欲望；也不至于在你心中，金钱取代了对丈夫的爱慕。\newline{}ewline{}因为观察人们的交往，我们常常凭经验发现，在某些人身上，当放纵受到抑制时，贪婪却增长了。\newline{}ewline{}因为，正如在身体的感觉中，盲人听觉更敏锐，且凭借触觉辨别许多事物，而有视力的人触觉并不如此灵敏；由此可知，当注意力的运用在某一方面（即视觉）受到限制时，它便转向其他感官，更敏锐地准备辨别，仿佛试图从一种感官中弥补另一种感官所欠缺的；同样，肉欲也常常在受制于感官交合的快乐之外，以更大的力量伸向对金钱的贪求，当它从某一方面被转移时，便以更炽热的激情转向另一方面。\newline{}ewline{}但在你身上，愿对财富的爱与对婚姻的爱一同冷却；愿你以虔诚的方式使用你的财产，将其导向属灵的喜乐，以致你的慷慨更热心于帮助贫困者，而不是使贪婪者富足。因为，送入天上宝库的，不是给贪婪者的礼物，而是给穷人的施舍，这施舍极大地助益寡妇的祈祷。\newline{}ewline{}此外，斋戒和守夜，只要不损害健康，若用于祈祷、咏唱圣咏、阅读和默想天主的法律，那么即使是看似劳苦的事，也会转变为属灵的喜乐。\newline{}ewline{}因为，对于喜爱之事，其劳作绝不令人疲乏，反而本身即是愉悦，如那些打猎、捕鸟、钓鱼、采摘葡萄、经商、或玩某种游戏自娱的人。\newline{}ewline{}因此，所爱之事至关重要。因为，对于所爱之事，要么没有劳苦，要么劳苦本身也被爱。试想，若在劳作中——如捕获野兽、装满酒桶和钱袋、投球游戏——有乐趣，而在赢得天主的劳作中却无乐趣，这将是多么可耻和可悲啊！\par

27. 的确，凡未婚妇女所享有的一切属神喜乐中，她们的圣洁行为也当谨慎；免得因傲慢而生活虽不邪恶，却因疏忽而声名狼藉。也不可听从那些圣洁男女的言语，当他们因某事的疏忽而受责，以致陷于恶意的猜疑，而他们明知自己的生活远非如此，却说良心在天主前已经足够，藐视众人对他们的看法——这不仅不智，而且残忍；因为他们杀害了他人的灵魂；无论是那些亵渎天主之道的人，他们循着猜疑，对圣人所贞洁的生活感到不悦，视之为可耻；还是那些为自己寻找借口，效法自己所想的而非所见的人。因此，凡保守自己生活免于可耻和邪恶行为的指责，是为自己行善；但同时也保守自己名誉的，也是向他人施行仁慈。因为对我们自己而言，我们的生命是必需的；对他人而言，我们的名誉是必需的；而且，我们仁慈地供给他人、为其益处所行的事，也必将增添我们自己的利益。因此，宗徒并非徒然说：\BibleQuote{我们预备善事}，他说，\BibleQuote{不仅在天主前，也在人前}；他也说：\BibleQuote{你们要凡事使众人喜悦，如同我也凡事使众人喜悦，不求自己的利益，只求众人的利益，使他们得救。}同样，在某劝勉中，他说：\BibleQuote{此外，弟兄们，凡是真实的，凡是圣洁的，凡是公义的，凡是纯洁的，凡是可爱的，凡是有美名的；若有什么德行，若有什么赞美，你们就该思念这些事，就是你们在我身上所学习、所领受、所听见、所看见的。}你们看，在他藉劝勉所叮嘱他们的许多事中，他也没有忽略提到\BibleQuote{凡是有美名的}；并用两个词概括了一切，他说：\BibleQuote{若有什么德行，若有什么赞美。}因为上述的善事属于德行，美名则属于赞美。我想宗徒并不以人的赞美为大事，却在别处说：\BibleQuote{但对我来说，受你们审断，或受人的日子审断，都是极小的事}；在另一处说：\BibleQuote{我若讨人的喜悦，就不是基督的仆人了}；又说：\BibleQuote{我们的光荣就是良心的见证。}但在这两者，即善行与美名，或更简洁地说，德行与赞美中，前者他为己持守，最为明智；后者他为人提供，最为仁慈。但人的谨慎无论多大，总不能完全避开最恶意的猜疑；当我们为自己的美名做了合理所能做的一切之后，若仍有人或凭虚假的恶言诬蔑我们，或因相信关于我们恶事而企图玷污我们的清誉，那么，良心的安慰便应临在，并且确然有喜乐，因为我们的赏报在天上是大大的，即使人们对我们说许多恶言，而我们仍圣洁公义地生活。因为那赏报是作为士兵的薪饷，藉着正义的武器，不仅在右边，也在左边；即透过光荣与卑贱，透过恶名与美名。\par

28. 因此，继续你的赛程，怀着坚忍奔跑，好使你能获得奖赏；并以生活的榜样和劝勉的言谈，尽力将所能影响的人一同带入这赛程。不要让虚浮之人的抱怨使你偏离这热忱的目标——他们藉此激励许多人跟随——他们说：如果人人都守贞，人类如何延续？好像这世界延续的原因，除了为使预定数量的圣人得以圆满外，再无其他；而若这数目早日圆满，世界的终结无疑不会延迟。也不要让人对你说：\Quote{虽然婚姻也是善的，但如果所有人都因赞美和热爱守贞而效法，那么在基督的身体内，如何能既有更大的善，也有较小的善呢？}——这不应阻止你劝勉别人也拥有你所拥有的善。首先，因为即使努力使人人守贞，实际上仍只有少数，因为\BibleQuote{不是人人都领受这话。}但经上既记载：\BibleQuote{谁能领受，就领受}，那么，当人不向那些不能领受的人保持沉默时，那些能领受的人便领受了。其次，我们也不该担心恐怕人人都会领受，以致基督身体中缺乏某种较小之善，即婚姻生活。因为如果人人都听见并领受，我们就当明白，这正是预定的，即婚姻之善在那些已经离开此生的肢体数目中已经足够。因为即使现在人人守贞，他们也不会将守贞的光荣归于那些已经为主仓房带来三十倍果实的人——如果那是指婚姻之善的话。因此，所有这各种善都将在那里有自己的位置，即使从今以后没有女子愿意结婚，没有男子愿意娶妻。所以，无忧无虑地敦促你所能影响的人，成为你现在的样子；并以警醒和热忱祈祷，藉至高者右手的援助和主最仁慈恩宠的丰沛，使你既能持守你现在的状态，也能进步到你将来的状态。\par

29. 其次，我藉着那一位恳求你，你从他那里既领受了这个恩赐，又盼望这恩赐的赏报，愿你记念将我放在连同你的整个家庭教会的祈祷中。确实，以最合宜的顺序，我应当给你现已年迈的母亲写一封关于祈祷的信；因为她主要负责通过祈祷为你奋斗，她对自己的关心少于对你的关心；并且我为你而非为她撰写这本关于守寡贞洁的小作品，因为你需要克服的，她的年龄已经克服了。但你的女儿，这位圣洁的童贞女，如果她因我们的劳苦而想了解关于她圣召的任何事，她有一本大书\WorkTitle{论圣童贞}可以阅读。关于阅读此书，我也曾劝告你，因为它包含了许多对两种贞洁（即童贞和守寡）都必要的内容；因此我在这里有的轻描淡写，有的完全略过，因为我在那里更充分地讨论了它们。\par

愿你持守于基督的恩宠之中。\par

\SourceRecord{Translated by C.L. Cornish.；Nicene and Post-Nicene Fathers, First Series；Vol. 3.；Edited by Philip Schaff.；Buffalo, NY: Christian Literature Publishing Co.,；1887.；<http://www.newadvent.org/fathers/1311.htm>.}
